% Part:sets-functions-relations
% Chapter: size-of-sets
% Section: pairing-alt

\documentclass[../../../include/open-logic-section]{subfiles}

\begin{document}

\olfileid{sfr}{siz}{pai-alt}

\olsection{اک متبادل جوڑی فنکشن}

\begin{explain}
$\Nat^2$ دیاں ہور وی گنتی وار فہرستاں نیں جیہناں دے اُلٹ فنکشن لبھنے
ودھ سوکھے نیں۔ ایہناں وچوں اک ایہ اے۔ گنتی وار فہرست نوں جدول وچ
ویکھن دی بجائے، مثبت صحیح عدداں دی اوس فہرست توں شروع کرو جیہدے نال
مُڈھ وچ خالی تھاواں جوڑیاں ہوئیاں نیں۔ تصور کرو کہ ایہ تھاواں
ترتیب وار جوڑیاں $\tuple{n,m}$ نال اگلے ڈھنگ نال بھریاں جا رہیاں نیں۔
پہلاں اوہ جوڑیاں لو جیہناں دے پہلے مقام اُتے~$0$ اے (یعنی جوڑیاں
$\tuple{0,m}$)؛ پہلی جوڑی (یعنی $\tuple{0,0}$) پہلی خالی تھاں اُتے
رکھو، فیر اک خالی تھاں چھڈ دیو، دوجی جوڑی (یعنی $\tuple{0,1}$) اگلی
خالی تھاں اُتے رکھو، فیر اک ہور خالی تھاں چھڈ دیو، تے ایویں اگے۔
% PNBSIZ-002 ماخذ دی درستی: ایتھے دوجی جوڑی (0,1) اے؛ انگریزی عبارت وچ غلطی نال (0,2) لکھیا سی، جدوں کہ نال دی پہلی قطار وی (0,1) ای وکھاؤندی اے۔
ساڈی گنتی وار فہرست دا (ادھورا) مُڈھ ہُن ایویں دِسدا اے:
\[\small
\begin{array}{@{}c c c c c c c c c c c@{}}
\mathbf 1 & \mathbf 2 & \mathbf 3 & \mathbf 4 & \mathbf 5 & \mathbf 6 & \mathbf 7 & \mathbf 8 & \mathbf 9 & \mathbf{10} & \dots \\ \\
\tuple{0,0} &  & \tuple{0,1} &  & \tuple{0,2} &  & \tuple{0,3} & & \tuple{0,4} &  & \dots \\
\end{array}
\]
ہُن جیہڑیاں تھاواں ہلے وی خالی نیں اوہناں لئی جوڑیاں $\tuple{1,m}$
نال ایہ کم دہراؤ، تے فیر ہر دوجی خالی تھاں چھڈدے جاؤ:
\[\small
\begin{array}{@{}c c c c c c c c c c c@{}}
\mathbf 1 & \mathbf 2 & \mathbf 3 & \mathbf 4 & \mathbf 5 & \mathbf 6 & \mathbf 7 & \mathbf 8 & \mathbf 9 & \mathbf{10} & \dots \\ \\
\tuple{0,0} & \tuple{1,0} & \tuple{0,1} &  & \tuple{0,2} & \tuple{1,1} &
\tuple{0,3} & & \tuple{0,4} &  \tuple{1,2} & \dots \\
\end{array}
\]
اوسے ڈھنگ نال جوڑیاں $\tuple{2,m}$، $\tuple{3,m}$ وغیرہ درج کرو۔
% OLSIZ-003 ماخذ دی درستی: لگاتار قطاراں (2,m)، (3,m) وغیرہ نیں؛ انگریزی عبارت وچ (2,m) دو واری آ گیا سی۔
ساڈی مکمل گنتی وار فہرست ایس طرح شروع ہوندی اے:
\[\small
\begin{array}{@{}cc c c c c c c c c c@{}}
\mathbf 1 & \mathbf 2 & \mathbf 3 & \mathbf 4 & \mathbf 5 & \mathbf 6 & \mathbf 7 & \mathbf 8 & \mathbf 9 & \mathbf{10} & \dots \\ \\
\tuple{0,0} & \tuple{1,0} & \tuple{0,1} & \tuple{2,0}  & \tuple{0,2} &
\tuple{1,1} & \tuple{0,3} & \tuple{3,0}  & \tuple{0,4} &  \tuple{1,2} & \dots \\
\end{array}
\]
جے اسی اُتلے جدول دے خانیاں نوں ایس گنتی وار فہرست مطابق نمبر دیئیے،
تاں سانوں صاف زگ زیگ لکیر نہیں لبھے گی، بلکہ ایہ ترتیب لبھے گی:
\[
\begin{array}{ c | c | c | c | c | c | c | c }
& \mathbf 0 & \mathbf 1 & \mathbf 2 & \mathbf 3 & \mathbf 4 & \mathbf 5 & \dots \\
\hline
\mathbf 0 & 1 & 3 & 5 & 7 & 9 & 11 & \dots \\
\hline
\mathbf 1 & 2 & 6 & 10 & 14 & 18 & \dots & \dots \\
\hline
\mathbf 2 & 4 & 12 & 20 & 28 & \dots & \dots & \dots \\
\hline
\mathbf 3 & 8 & 24 & 40 & \dots & \dots & \dots & \dots \\
\hline
\mathbf 4 & 16 & 48 & \dots & \dots & \dots & \dots & \dots \\
\hline
\mathbf 5 & 32 & \dots & \dots & \dots & \dots & \dots & \dots \\
\hline
\vdots & \vdots & \vdots & \vdots & \vdots & \vdots & \vdots & \ddots\\
\end{array}
\]
اسی ویکھ سکدے آں کہ قطار~$0$ دیاں جوڑیاں ساڈی گنتی وار فہرست دے
طاق نمبر والے مقاماں اُتے نیں، یعنی جوڑی $\tuple{0,m}$، مقام $2m+1$
اُتے اے؛ دوجی قطار دیاں جوڑیاں $\tuple{1,m}$ اوہناں مقاماں اُتے نیں
جیہناں دا نمبر کسے طاق عدد دا دُگنا اے، خاص طور تے $2 \cdot (2m+1)$؛
تیجی قطار دیاں جوڑیاں $\tuple{2,m}$ اوہناں مقاماں اُتے نیں جیہناں دا
نمبر کسے طاق عدد دا چوگنا اے، $4 \cdot (2m+1)$؛ تے ایویں اگے۔ ہر قطار
لئی $(2m+1)$ دے عامل $1$، $2$، $4$، $8$، \dots، ٹھیک~$2$ دیاں
قوتاں نیں: $1= 2^0$، $2 = 2^1$، $4 = 2^2$، $8 = 2^3$، \dots\@
اصل وچ متعلقہ قوت نما ہمیشہ زیرِ غور جوڑی دا پہلا رکن اے۔ ایس لئی جوڑی
$\tuple{n,m}$ لئی عامل $2^n$ اے۔ ایس توں سانوں عام فارمولا ملدا اے:
$2^n \cdot (2m+1)$۔ پر ایہ جوڑیاں توں \emph{مثبت} صحیح عدداں ول نقشہ
بندی اے، یعنی $\tuple{0,0}$ دا مقام~$1$ اے۔ جے اسی مقام~$0$ توں شروع
کرنا چاہندے آں تاں نتیجے وچوں~$1$ گھٹانا پئے گا۔ ایس توں سانوں ایہ
ملدا اے:
\end{explain}

\begin{ex}
اگلے فارمولے نال دتا گیا فنکشن $h\colon \Nat^2 \to \Nat$
\[
h(n,m) = 2^n (2m+1) - 1
\]
قدرتی عدداں دیاں جوڑیاں دے سیٹ~$\Nat^2$ لئی اک جوڑی فنکشن اے۔
\end{ex}

\begin{explain}
ایس مطابق، $\Nat^2$ دی ساڈی دوجی گنتی وار فہرست وچ جوڑی
$\tuple{0,0}$ دا کوڈ $h(0,0) = 2^0(2\cdot 0+1) - 1 = 0$ اے؛
$\tuple{1,2}$ دا کوڈ $2^{1} \cdot (2 \cdot 2 + 1) - 1 = 2
\cdot 5 - 1 = 9$ اے؛ تے $\tuple{2,6}$ دا کوڈ $2^{2} \cdot (2
\cdot 6 + 1) - 1 = 51$ اے۔
\end{explain}

کدی کدی قدرتی عدداں دیاں جوڑیاں $\Nat^2$ نوں کوڈ کرنا ای کافی ہوندا
اے، تے ایہ شرط ضروری نہیں ہوندی کہ کوڈ بندی ہر ہدف قدر لیندی ہووے۔
ایہو جہیاں کوڈ بندیاں دے اُلٹ صرف جزوی فنکشن ہوندے نیں۔

\begin{ex}
اگلے فارمولے نال دتا گیا فنکشن $j\colon \Nat^2 \to \Nat^+$
\[
j(n,m) = 2^n3^m
\]
اک !!a{injective} فنکشن $\Nat^2 \to \Nat$ اے۔
\end{ex}

\end{document}
